% Research integration fragment, 2026-10-08.  Old sources are unchanged.
\subsection{Finite same-level feedback in the repaired second seed}
\label{sec:second-seed-finite-feedback}

This gives the full positive feedback operator for the non-switched
terms in Wu 2004, Lemma 4.2, after the repair of $\Gamma_5$.
It uses already proved child estimates at their actual states.
The numerical probe below is only a floating feasibility calculation.

\paragraph{States and parameter restrictions.}
Use the stable enhanced or classical family of
Proposition~\ref{prop:second-seed-stable-prefix}.  At a parent state
$\mathfrak s$ write
\[
 B=N^{T+o(1)},\quad Q=N^\nu,\quad R=Q/B=N^{L+o(1)},\quad
 L=\nu-T>0,\quad \alpha=(1-T)/L\ge2.
\]
In the enhanced case retain full specified well support and one
proved factorization branch, with its stable ceiling $\chi$ and
every old gate.  In particular any third-slot bound $c$ used before
three actual specified slots exist requires $\chi\le c$.
The classical case uses the absolute remainder below level $1/2$.
All statements are averaged on finite positive marked boxes.

Choose parameters
\[
 2<s<\phi_3<\phi_2<\phi_1<s_0\le5,
 \quad a_0=1/s_0,\quad a_i=1/\phi_i,\quad b=1/s,
\]
with
\begin{equation}\label{eq:feedback-parameter-gates}
 a_3+2a_1<1,\qquad 3a_1<1,\qquad a_2+2a_1<1.
\end{equation}
The practical table below uses $2.1\le s\le2.5$.
Require the parent's comparison cutoff $L/s<\chi$ on an enhanced
box; all its new specified primes are then below the old prefix.
The state at a child always retains the same total level $Q$.

Let $H_{\mathfrak t}(w)$ and $h_{\mathfrak t}(w)$ mean previously
proved defects, at a state $\mathfrak t$, in the bounds
\[
 \Phi_{\mathfrak t,w}\le
 (A(w)-H_{\mathfrak t}(w)+o(1))\Theta_{\mathfrak t},\qquad
 \Phi_{\mathfrak t,w}\ge
 (a(w)+h_{\mathfrak t}(w)+o(1))\Theta_{\mathfrak t}.
\]
Here $A(w)=wF(w)/(2e^\gamma)$, $a(w)=wf(w)/(2e^\gamma)$,
and unproved defects are zero.  On $1\le w<2$ a positive upper
defect must come from the already justified cutoff comparison;
it is not an invocation of an ordinary upper-sieve theorem at a
forbidden parameter.  The normalized comparison multiplier remains
$A(w)=1$ on $1\le w\le3$.  We use $h=0$ below two.

\paragraph{The exact child states and cutoffs.}
The two terms in $\Gamma_1$ keep the parent prefix and ratio
$\alpha$, with respective arguments $s_0,\phi_1$.
For each of $\Gamma_2,\Gamma_3,\Gamma_4$, append a prime
$q=N^{Lu}$, where respectively
\[
 a_0<u<b,\qquad a_0<u<a_2,\qquad a_0<u<a_3.
\]
Call its state $\mathfrak s_u$.  Its data are
\begin{equation}\label{eq:feedback-single-child}
 L_u=L(1-u),\qquad
 \alpha_u=\frac{\alpha-u}{1-u},\qquad
 \rho_u=s_0(1-u),\qquad
 \log_N(\hbox{physical cutoff})=La_0.
\end{equation}
The singleton square condition is $2u<1$, and the physical cutoff
is at most the new smallest specified prime $N^{Lu}$.

For the pair terms use $q_1=N^{Lu}<q_2=N^{Lv}$ and put
\[
 D_5=\{a_0<u<v<a_2\},\quad
 D_6=(a_0,a_1)\times(a_2,a_3),\quad
 D_7=\{a_0<u<v<a_1\},\quad
 D_8=(a_0,a_1)\times(a_1,a_2).
\]
The child state $\mathfrak s_{u,v}$ has
\begin{equation}\label{eq:feedback-pair-child}
 L_{u,v}=L(1-u-v),\qquad
 \alpha_{u,v}=\frac{\alpha-u-v}{1-u-v}.
\end{equation}
Its actual physical cutoff and argument are
\[
\begin{array}{c|c|c}
 i&\log_N(\hbox{physical cutoff})&\rho_i(u,v)\\\hline
 5,6&La_0&s_0(1-u-v)\\
 7,8&Lu&(1-u-v)/u.
\end{array}
\]
In particular $\Gamma_7,\Gamma_8$ do not have cutoff $R^{a_0}$.
Changing that cutoff would give the wrong child argument and kernel.

The residual pair is well supported exactly under
\begin{equation}\label{eq:feedback-pair-squares}
 2v<1,\qquad v+2u<1.
\end{equation}
Both hold everywhere in $D_6,D_7,D_8$ by
\eqref{eq:feedback-parameter-gates}.  On $D_5$ retain only
$D_5^+=D_5\cap\{v+2u<1\}$ for ordinary feedback; write
$D_5^-=D_5\setminus D_5^+$.  The first test $2v<1$ holds throughout.
There is no replacement of these asymmetric pair tests by a
condition such as three times one exponent being below the level.

On every retained pair domain, $1<\rho_i\le s_0-2\le3$.
Indeed $1-u-v>u\ge a_0$ proves the lower endpoint; maximizing
the displayed arguments at the lower endpoints gives the upper
endpoint.  For $\rho_i<2$, compare the positive child count at
any previously licensed larger argument, hence a smaller cutoff.
This keeps every compared prime below its actual cutoff and below
the child's smallest specified prime.  The argument $\rho_i$ in
the operator is the actual argument, with that proved comparison
already included in $H_{\mathfrak s_{u,v}}$.

The complete prefix tests follow from the merge lemma.  More
explicitly, all old slots precede the one or two new slots, so
old complete and selected-remaining tests are unchanged.  The
new tests are precisely $2u<1$ or
\eqref{eq:feedback-pair-squares}, multiplied by $B$.
Ordinary residual atoms are merged at the indicated child level.
Original selected slots and the relevant first three $W$ slots,
or their stable upper bound, are retained.  All comparisons lower
cutoffs, and no new slot can cross back above an old slot.
Finally
\[
 \alpha_u=1+\frac{\alpha-1}{1-u}\ge\alpha,\qquad
 \alpha_{u,v}=1+\frac{\alpha-1}{1-u-v}\ge\alpha.
\]
Thus ratios remain at least two, but are not reset to two.
Signed atom normalizations, repeated-label factorials, coprime
means and moving-residue margins are those of the existing stable
interface.  If a more specialized child defect has further gates,
it is used only where those gates also hold; elsewhere it is zero.

\paragraph{A uniform rough-number envelope with an exact check.}
Put $c_\omega=70893/125000=0.567144$.  Then
\begin{equation}\label{eq:feedback-omega-envelope}
 \omega(x)\le\max\{1/x,c_\omega\}\qquad(x\ge1).
\end{equation}
For $1\le x\le2$ this follows from $\omega(x)=1/x$.
For $2\le x\le3$, write $v=x-1$ and maximize
$1+\log v-c_\omega(v+1)$ at $v=1/c_\omega$.
The maximum is $-\log c_\omega-c_\omega<0$, because the exact
positive rational calculation
\[
 c_\omega\sum_{k=0}^{8}\frac{c_\omega^k}{k!}-1
 =\frac{65111700737847638377192824351410078164698277}
 {33378601074218750000000000000000000000000000000000}>0
\]
implies $c_\omega e^{c_\omega}>1$.
For $x\ge3$, integrate $(x\omega(x))'=\omega(x-1)$ and induct
over unit intervals, starting from $\omega(3)<c_\omega$.
This proves the envelope without a floating maximization.

On $D_5^-$ the rough argument is at least
$s_0(2-2a_2)>2$, so the retained positive replacement kernel may
be sharpened from $3Ks_0/20$ to $Kc_\omega s_0/4$.
There is still no child feedback on this domain.
For the switched triple domains $9,\ldots,15$, the following
numbers bound their argument from below for all $\alpha\ge2$:
\[
\begin{array}{c|l}
9&2\phi_3-3\\
10&2\phi_2-2-\phi_2/s\\
11&2\phi_3-\phi_3/\phi_2-2\\
12&2\phi_1-2-\phi_1/s\\
13&2\phi_2-\phi_2/\phi_1-1-\phi_2/s\\
14&2s-s/\phi_1-2\\
15&2\phi_3-\phi_3/\phi_2-1-\phi_3/s.
\end{array}
\]
If the entry is $m_i$, replace the former factor $3/5$ in its
explicit logarithmic mass by
$C_i=\max\{c_\omega,1/\max(1,m_i)\}$.
This also covers an enlarged parameter set where $m_i<5/3$;
nonunit terms only exist at argument at least one.
Use the previous coarse $\omega\le1$ bounds for $16,\ldots,21$,
including the separate uniform unit terms in $M_{20},M_{21}$.
There are no unit terms for $9,\ldots,19$, since $s>2$ gives
every specified sum at most $4/s<2\le\alpha$.

\paragraph{The positive feedback operator.}
Let $J_2^0(s,s_0,\phi_1,\phi_2,\phi_3)$ denote the ordinary-only
repaired functional, with the preceding uniform switched envelopes.
It is the functional of the saved seed with $A(s)$ replacing one
if needed.  Its ordinary pair kernels are $1/(1-u-v)$ on
$D_5^+,D_6,D_7,D_8$, and its bad-$\Gamma_5$ kernel is
$Kc_\omega s_0/4$.  No switched term contains an imported defect.

Define the same-level positive linear operator
\begin{align}
 \mathcal P_{2,\mathfrak s}[H,h]
 =\frac15\bigg\{&4H_{\mathfrak s}(s_0)+H_{\mathfrak s}(\phi_1)
 +\sum_{c\in\{b,a_2,a_3\}}\int_{a_0}^{c}
 \frac{h_{\mathfrak s_u}(s_0(1-u))}{u(1-u)}\,du\notag\\
 &+\iint_{D_5^+}
 \frac{H_{\mathfrak s_{u,v}}(s_0(1-u-v))}
      {uv(1-u-v)}\,du\,dv\notag\\
 &+\sum_{i=6}^8\iint_{D_i}
 \frac{H_{\mathfrak s_{u,v}}(\rho_i(u,v))}
      {uv(1-u-v)}\,du\,dv\bigg\}.
 \label{eq:second-seed-feedback-operator}
\end{align}
The valid new upper defect at the parent is
\begin{equation}\label{eq:second-seed-feedback-bound}
 H_{\mathfrak s}(s)\ge
 \max\{0,J_2^0+\mathcal P_{2,\mathfrak s}[H,h]\}.
\end{equation}
Here the defects on the right are previously proved inputs, so the
formula describes a finite additional proof step and is not circular.
To prove it, insert each child bound with its actual residual level
into the unchanged 21-term pointwise inequality.  An improvement
to $\Gamma_1$ subtracts $4H(s_0)+H(\phi_1)$ from its upper bound;
an improvement to a negative term adds its lower defect; and an
ordinary pair upper improvement subtracts its upper defect.
All three effects improve the final upper bound.  The change from
child normalization to parent normalization contributes
$1/(1-u)$ or $1/(1-u-v)$, exactly as displayed.  Division by five
then proves the formula.  Every coefficient of the operator is
nonnegative.

The implementation shared with the first-seed iteration uses the
special case of defects $H(w),h(w)$ uniform over all these stable
and classical states and every $\alpha\ge2$.  This removes the
state subscript in \eqref{eq:second-seed-feedback-operator}; it does
not change the actual ratios in
\eqref{eq:feedback-single-child}--\eqref{eq:feedback-pair-child}.
For a finite parameter bank, take the maximum of its candidate
bounds and already proved alternatives on a common positive
partition.  Different upper candidates are not added.

\paragraph{Finite parameter probe and its limits.}
The new module \texttt{probe\_second\_seed\_feedback.py} exports
\texttt{SecondSeedOperator(parameters,K,n).candidate(H,h)} and
\texttt{.value(H,h)} for numpy-compatible uniform input functions.
It treats the bad triangle by its exact two-piece boundary, with
no masked-grid replacement for that boundary.  It checks the
three gates in \eqref{eq:feedback-parameter-gates} before admitting
a parameter choice.

The initial upper inputs are the proved values $H=9/500$ on
$[1,11/5]$, $H=7/4625$ on $(11/5,3]$, and zero elsewhere.
The lower input is one Buchstab integral of their pointwise maximum,
with the lower endpoint $s-1$, including $[1,2)$ children.
There are 881 admitted parameter choices in the saved finite search.
It refines selected choices with quadrature orders 24, 48 and 80.
For $K=6.81$, the order-80 outcomes at the selected candidates are
\[
\begin{array}{c|c|c|c}
s&\hbox{ordinary-only seed}&\hbox{feedback}&\hbox{candidate}\\\hline
2.1&.02047094&.00253880&.02300973\\
2.2&.02057652&.00243968&.02301620\\
2.3&.01879876&.00245380&.02125256\\
2.4&.01571500&.00244468&.01815968\\
2.5&.01164637&.00229828&.01394466.
\end{array}
\]
These are sampled choices, not certified optima.
At the original parameters $(2.2,4.54,3.53,2.90,2.44)$ the
diagnostic is $.02033936+.00258581=.02292517$.
The feedback splits approximately as $.00093624$ from the three
negative terms and $.00164957$ from the four ordinary pair terms;
the first-generation $\Gamma_1$ contribution is zero because these
initial $H$ inputs vanish above three.  The same original parameters
give $.02299950$ if the separately supplied $K=6.792$ is used.
No validity or certification of that alternative $K$ is inferred
from this probe.

The discontinuities in chosen input defects make these quadratures
especially unsuitable as outward bounds without further partitioning.
Their agreement is a diagnostic, not a certificate.  The operator
shows a nonzero additional feedback of order $0.0025$ in a seed,
but that is not itself a gain of $0.0025$ in the complete Chen count.
The contribution must still pass through the full eleven-term
integrals and their existing minima.  These computations alone do
not fill the remaining frontier gap, establish a new global count,
or justify importing any old double-sieve table.  A subsequent
finite typed iteration may use this operator, but its resulting
integrals require their own certificates.
